\documentclass[12pt]{article}
\usepackage[utf8]{inputenc}
\usepackage[T1]{fontenc}
\usepackage{hyperref}
\title{Punu Language Resources for K-12 Curriculum Development}
\author{Compiled by Claude Sonnet 5.5 (claude-sonnet-5-5) for Andr\'as Kornai}
\date{}
\begin{document}
\maketitle

\section{Introduction}
Punu (Yipunu, ISO 639-3 \texttt{puu}) is a Bantu language (Guthrie B.43) of Gabon and the Republic of Congo. Wikipedia gives about 120 thousand native speakers, mainly in southern Gabon. The searches found academic work on Yipunu lexicography but no school materials.

\section{Primary Educational Resources}
\begin{itemize}
\item A Stellenbosch thesis proposing a model for a bilingual Yipunu--French school dictionary, with Yipunu as the dominant language (a design model, not a published dictionary): \url{https://hdl.handle.net/10019.1/1416}
\end{itemize}

\section{Linguistic and Reference Materials}
\begin{itemize}
\item Wikipedia, Punu language: \url{https://en.wikipedia.org/wiki/Punu_language}
\item Metalexicographic criteria for a monolingual descriptive dictionary of standard Yipunu (Stellenbosch): \url{https://scholar.sun.ac.za/items/f97e2537-5d0a-45e7-985e-2eb5095285d6}
\item Lexikos article discussing the CMA \emph{Dictionnaire fran\c{c}ais--yipounou/yipounou--fran\c{c}ais}: \url{https://lexikos.journals.ac.za/pub/article/view/651/251}
\item Second Lexikos article from the same search: \url{https://lexikos.journals.ac.za/pub/article/view/652/252}
\item LINGUIST List issue 17.3333 entry: \url{https://linguistlist.org/issues/17/3333}
\end{itemize}
No dedicated grammar was found.

\section{Religious and Community Resources}
None found in these searches.

\section{Digital and Audio Resources}
None found in these searches.

\section{Curriculum Development Context}
Existing Gabonese-language dictionaries are mostly translation dictionaries biased toward French, compiled by missionaries and colonial administrators; there is no monolingual Yipunu dictionary. A curriculum project would need to build primers and readers largely from scratch, possibly starting from the CMA bilingual dictionary and the Stellenbosch design work.

\end{document}